% Test Automation Plan (concise edition): careers.bupa.com.sa
% Build: pdflatex test-automation-plan-concise.tex (twice, for cross-references).
\documentclass[11pt,a4paper]{article}

\usepackage[T1]{fontenc}
\usepackage[utf8]{inputenc}
\usepackage{lmodern}
\usepackage[a4paper,margin=2.2cm,headheight=14pt]{geometry}
\usepackage{microtype}
\usepackage{xcolor}
\usepackage{array,booktabs,tabularx}
\usepackage{enumitem}
\usepackage{fancyhdr}
\usepackage{titlesec}
\usepackage{listings}
\usepackage{hyperref}

\definecolor{accent}{HTML}{1F4E79}
\definecolor{muted}{HTML}{595959}
\definecolor{codebg}{HTML}{F4F4F4}

\hypersetup{colorlinks=true,linkcolor=accent,urlcolor=accent,
  pdftitle={Test Automation Plan: careers.bupa.com.sa},pdfauthor={Hani Fayed}}

\titleformat{\section}{\Large\bfseries\color{accent}}{\thesection.}{0.5em}{}
\titleformat{\subsection}{\large\bfseries\color{accent}}{\thesubsection}{0.5em}{}
\titlespacing*{\section}{0pt}{1.3em}{0.5em}
\titlespacing*{\subsection}{0pt}{0.9em}{0.3em}

\setlist{itemsep=0.1em,topsep=0.25em}
\setlength{\parindent}{0pt}
\setlength{\parskip}{0.45em}
\setlength{\emergencystretch}{3em}
\renewcommand{\arraystretch}{1.22}

\pagestyle{fancy}
\fancyhf{}
\fancyhead[L]{\small\color{muted}Test Automation Plan}
\fancyhead[R]{\small\color{muted}careers.bupa.com.sa}
\fancyfoot[L]{\small\color{muted}Hani Fayed}
\fancyfoot[R]{\small\color{muted}\thepage}
\renewcommand{\headrulewidth}{0.4pt}

\lstset{basicstyle=\ttfamily\small,backgroundcolor=\color{codebg},frame=single,
  rulecolor=\color{codebg},columns=fullflexible,keepspaces=true,xleftmargin=4pt,xrightmargin=4pt}

\newcolumntype{Y}{>{\raggedright\arraybackslash}X}
\newcolumntype{L}[1]{>{\raggedright\arraybackslash}p{#1}}

\begin{document}

% ---------- cover page ----------
\begin{titlepage}
\vspace*{4cm}
{\color{accent}\Huge\bfseries Test Automation Plan\par}
\vspace{0.6cm}
{\Large\color{muted} Bupa Arabia careers portal\par}
\vspace{0.2cm}
{\large\color{muted} careers.bupa.com.sa\par}
\vspace{0.8cm}
{\color{accent}\rule{\linewidth}{0.8pt}}
\vspace{0.6cm}

\begin{tabular}{@{}ll@{}}
\textbf{Author} & Hani Fayed \\
\textbf{Context} & Talentera technical assessment, Part 2 \\
\textbf{Date} & September 2026 \\
\end{tabular}

\vfill
{\small\color{muted} This plan describes how I would approach automated testing of the Bupa Arabia careers
portal: what to automate, with which tools, in which order, and how the suite would be run and maintained.\par}
\end{titlepage}

% ---------- contents ----------
\tableofcontents
\thispagestyle{empty}
\newpage
\setcounter{page}{1}

% ======================================================================
\section{Objective}

\textbf{Purpose.} The careers portal is where candidates discover jobs at Bupa Arabia and apply for them.
If a candidate cannot register, sign in, find a job or submit an application, Bupa loses that applicant.
The purpose of automation is to protect these journeys and give the team fast, trustworthy feedback on
every change.

\textbf{Main goals}
\begin{itemize}
  \item Automate the critical candidate journeys end to end, starting with the ``Apply to a Job'' flow I built in Part 1.
  \item Check business rules (search results, saved jobs, alerts, access control) quickly at the API level.
  \item Run a short smoke suite on every change and a fuller regression suite every night.
  \item Keep the suite stable and easy to maintain as the portal changes.
\end{itemize}

\textbf{Expected benefits:} faster releases with less manual regression; defects caught before candidates see
them; consistent coverage across browsers; and reports with screenshots and traces that make failures
quick to diagnose.

% ======================================================================
\section{Automation Scope}

\textbf{Main features to cover:} registration with CV upload and parsing; sign-in with the emailed
verification code; forgot password; job search, advanced search and job details; applying for a job;
My Applications and withdrawing; Saved Jobs; Job Alerts; the CV and profile pages; the English and Arabic
versions of the site.

\textbf{Critical user flows}
\begin{enumerate}
  \item \textbf{Apply to a Job:} sign in $\rightarrow$ search $\rightarrow$ open a job $\rightarrow$ apply
        $\rightarrow$ confirm it appears in My Applications $\rightarrow$ withdraw.
  \item \textbf{Registration:} upload a CV $\rightarrow$ review the parsed details $\rightarrow$ create the account
        $\rightarrow$ complete the missing fields.
  \item \textbf{Sign-in:} email and password $\rightarrow$ verification code $\rightarrow$ signed in.
  \item \textbf{Forgot password:} request a reset $\rightarrow$ receive the email $\rightarrow$ set a new password
        $\rightarrow$ sign in with it.
  \item \textbf{Job discovery:} keyword search, filters, sorting, paging, advanced search, job details.
  \item \textbf{Saved Jobs:} save from a job card or job page, see it in Saved Jobs, unsave it.
  \item \textbf{Job Alerts:} create an alert from a search $\rightarrow$ edit it $\rightarrow$ receive the alert
        email with matching jobs $\rightarrow$ delete it.
\end{enumerate}

\subsection*{What my tour showed}
I explored the portal as a candidate, in English and Arabic, on desktop and mobile. These observations
shaped the scope and priorities of this plan:

\begin{tabularx}{\linewidth}{@{}L{3.3cm}Y@{}}
\toprule
\textbf{Area} & \textbf{What I observed} \\
\midrule
CV parsing & From my PDF CV, name, email, date of birth, job titles, dates, education, work descriptions and
languages were extracted correctly, and skills were added from keywords. Phone, job location, city and company
industry were left empty. Parsing took 1--2 minutes, and there is no review step: the account is created and I was
then sent to the profile to complete the missing fields. With a random \texttt{.txt} file, the page still said
``We were able to extract the details from your CV'', while asking me to enter my email because it could not be
read: a misleading success message. Other file types are rejected with a clear message listing
\texttt{.PDF}, \texttt{.DOC}, \texttt{.DOCX} and \texttt{.TXT}. \\
Sign-in & Every login asks for a reCAPTCHA and a 4-digit code sent by email. The code arrived within about a minute
every time and stays valid for 5 minutes. Resending sends a new code and invalidates the old one, and an old code
shows ``You have entered a wrong code, please enter the correct code and try again.'' Only one session per account
can be active, and it is tied to the browser. \\
Accessibility and language & Registration and forgot-password fields have no labels, only placeholders; automated
accessibility checks found low contrast and unnamed links on every page I tested; the Arabic header still says
``Login'' in English. \\
Saved Jobs & Works as expected: the save star appears on job cards and job pages when signed in (and is hidden when
signed out), the saved state shows correctly, and unsaving works. \\
Search paging & Paging past the last page is handled inconsistently: with 5 pages of results, page 6 shows
``Sorry, no jobs matched your search criteria'', but page 99 silently shows page 1. \\
\bottomrule
\end{tabularx}

\textbf{What I would automate:} the critical flows above on the UI; search, advanced search, job details, saved
jobs, alerts and access control on the API; negative cases for every form; basic accessibility checks on key pages.

\textbf{What would remain manual}
\begin{itemize}
  \item \textbf{The reCAPTCHA on production} (at sign-in and on forgot password): it is never automated or bypassed. In a QA environment it would use
        Google's test keys; on production a person signs in once and the saved session is reused (as in Part 1).
  \item \textbf{Exploratory testing} of new features, and usability judgement.
  \item \textbf{Content and language quality}, especially Arabic wording, which needs a native reader.
  \item \textbf{Email content and layout}, reviewed by eye once per release.
  \item \textbf{Screening questionnaires on production:} answers would reach a real recruiter.
\end{itemize}

% ======================================================================
\section{Testing Approach}

\begin{tabularx}{\linewidth}{@{}L{2.7cm}Y@{}}
\toprule
\textbf{Type} & \textbf{How I would apply it} \\
\midrule
UI testing & Playwright tests that drive the portal like a candidate, built on the Page Object Model, limited to the
flows that only a browser can prove. \\
API testing & Postman collections for search, job details, saved jobs, alerts and access control. Faster and more
stable than the UI, so most business rules are checked here. Advanced search has many fields (keyword match mode,
location, role, category, posted date, salary range and currency, employment type and more), so its combinations
are covered by data-driven Postman runs; the UI only checks that the form submits, results match and Reset clears it. \\
Smoke testing & A small, fast set: home page loads, search returns jobs, a job page opens, sign-in works. Runs on
every pull request and after every deployment. \\
Regression testing & All automated P0 and P1 tests, on all browsers, every night and before each release. \\
Negative testing & Wrong password and empty fields; wrong, old and expired verification codes; forgot password
with an unknown email and an expired reset link; unsupported and unreadable CV files (including the misleading
\texttt{.txt} message); the save star hidden when signed out; protected pages redirecting to login; unusual search
input; advanced search with a minimum salary above the maximum, or criteria that match nothing; result pages past
the last page. \\
\bottomrule
\end{tabularx}

% ======================================================================
\section{Tools \& Technologies}

\begin{tabularx}{\linewidth}{@{}L{3.6cm}Y@{}}
\toprule
\textbf{Tool} & \textbf{Why} \\
\midrule
Playwright & Automatic waiting, Chromium, Firefox and WebKit from one API, saved sign-in sessions, and a trace
viewer for debugging. I already used it successfully on this portal in Part 1. \\
TypeScript & Type checking catches mistakes before a test runs, and it is the language Playwright is designed around. \\
Postman (Newman in CI) & Readable API collections the whole team can open and run, versioned in the repository and
run from the command line in CI. Playwright's \texttt{request} fixture is used strictly for UI test setup and
teardown (for example, removing a saved job after a test), so the UI suite never depends on Postman. \\
Git & Version control, with every test change reviewed through a pull request. \\
GitHub Actions & Runs the suites automatically on pull requests, on a nightly schedule and after deployments,
with reports kept as artifacts. \\
HTML and Allure reports & Playwright's HTML report is produced on every run and includes screenshots, videos and
traces. Allure is added to the nightly regression run for history and trends across runs. \\
\bottomrule
\end{tabularx}

% ======================================================================
\section{Test Prioritization}

I rate each area by \textbf{business impact}, \textbf{risk of failure} and \textbf{how often it is used}, and
automate the highest first.

\begin{tabularx}{\linewidth}{@{}l L{4.2cm} Y@{}}
\toprule
\textbf{Priority} & \textbf{Area} & \textbf{Reason} \\
\midrule
P0 & Critical business flows & \textbf{Apply to a Job} (already automated in Part 1), sign-in, registration, job
search. If these fail, candidates cannot apply. \\
P0 & High-risk functionality & CV parsing (slow, some fields left empty, misleading message for unreadable files)
and the emailed verification code, which every sign-in depends on. \\
P1 & Frequently used functionality & Job details, filters and sorting, Advanced Search, Saved Jobs, Job Alerts,
My Applications, forgot password. \\
P1 & Repetitive regression tests & Form validation, verification-code handling, access to protected pages,
accessibility and Arabic checks: tedious to repeat by hand every release. \\
P2 & Lower-priority scenarios & Static content pages, FAQ, footer links, individual advanced search options that are rarely used. \\
\bottomrule
\end{tabularx}

% ======================================================================
\section{Automation Framework Structure}

The framework extends the one I built in Part 1:
\begin{itemize}
  \item \textbf{Page Objects:} one class per page, holding its locators and actions. Locators use roles, labels and
        placeholders first, never styling classes.
  \item \textbf{Test cases:} specs grouped by journey; page objects wait for the page, specs make the assertions.
  \item \textbf{Test data:} JSON files with no secrets, plus sample CVs (valid and invalid).
  \item \textbf{Reusable utilities:} environment settings, email mailbox access, accessibility helper.
  \item \textbf{Configuration:} \texttt{playwright.config.ts} (browsers, base URL, retries) and a \texttt{.env} file for credentials.
  \item \textbf{Test reports:} generated into their own folders and never committed.
\end{itemize}

\begin{lstlisting}
tests/
  pages/          LoginPage.ts, JobSearchPage.ts, ApplicationPage.ts ...
  specs/          apply-job.spec.ts, register.spec.ts, advanced-search.spec.ts ...
  test-data/      jobs.json, invalid-users.json, cvs/
  fixtures/       test.ts (shared page objects and accounts)
  utils/          env.ts, mailbox.ts
postman/          collections/, environments/
playwright.config.ts
.github/workflows/
\end{lstlisting}

% ======================================================================
\section{Test Execution}

\textbf{Local execution:} \texttt{npx playwright test} runs everything; single specs, tags (\texttt{--grep @smoke})
and headed mode make debugging easy; Postman collections run from the app or with Newman.

\textbf{Environment and test data:} the full suite runs in a QA environment with dedicated test accounts and
test job postings. Production gets only read-only smoke tests, because a real application reaches Bupa's HR team.
For the same reason, my Part 1 suite applies and then withdraws by default, and has a dry-run mode that stops
before submitting.

\textbf{CI execution:} GitHub Actions runs the suites on every pull request, every night and after each deployment.
The jobs run inside Docker images (the official Playwright and Newman images) so CI and local runs behave the same.

\begin{tabularx}{\linewidth}{@{}L{2.3cm}Y L{3.3cm}@{}}
\toprule
\textbf{Suite} & \textbf{Content} & \textbf{When} \\
\midrule
Smoke & A handful of \texttt{@smoke} UI tests plus the API collections; a few minutes & Every pull request, after each deployment \\
Regression & All P0 and P1 tests on every browser, accessibility checks & Nightly, before each release \\
\bottomrule
\end{tabularx}

\textbf{Browser coverage:} Chromium, Firefox and WebKit on desktop, plus a mobile viewport for the critical flows.

\textbf{Parallel execution:} read-only tests (search, job pages, API) run in parallel freely. Signed-in tests need
one test account per parallel worker, because the portal allows only one active session per account.

% ======================================================================
\section{Reporting \& Maintenance}

\begin{itemize}
  \item \textbf{Test reports:} the HTML report on every run, Allure for nightly history, and Newman results for the
        API; all published as CI artifacts.
  \item \textbf{Screenshots and traces on failure:} Playwright keeps a screenshot, a video and a full trace (every
        step, network call and console message) for each failed test.
  \item \textbf{Logging:} named test steps, and the job or account used by each run recorded in the report.
  \item \textbf{Flaky tests:} they are investigated at once, not hidden by retries. Read-only tests may retry once and
        are reported as flaky if a retry was needed; tests with side effects (such as applying) never retry. A test
        that stays flaky is quarantined and fixed within the sprint.
  \item \textbf{When the website changes:} page objects keep each locator in one place, so a change is fixed once.
        I would also ask developers for stable \texttt{data-testid} attributes on key controls.
  \item \textbf{Code review and maintenance:} every test change goes through a pull request with type checking and
        review; slow or duplicated tests are cleaned up regularly.
\end{itemize}

% ======================================================================
\section{Implementation Plan}

\begin{tabularx}{\linewidth}{@{}L{4.3cm}Y@{}}
\toprule
\textbf{Phase} & \textbf{Deliverables} \\
\midrule
1. Set up framework & Project structure, configuration, page object base, reporting, test accounts
(largely in place from Part 1). \\
2. Automate critical flows & Apply to a Job (done in Part 1), sign-in, registration with CV upload, job search. \\
3. Add API and regression coverage & Postman collections for search, Advanced Search, job details, saved jobs and
alerts; negative tests; UI tests for Saved Jobs, Job Alerts, My Applications and forgot password. \\
4. Integrate with CI/CD & GitHub Actions workflows for smoke and regression, Docker images, reports as artifacts. \\
5. Expand and maintain & Arabic and mobile coverage, accessibility checks, regular review of flaky and
slow tests, a test added with each new feature. \\
\bottomrule
\end{tabularx}

% ======================================================================
\section{Conclusion}

My goal is not to automate everything, but to build reliable, maintainable automation around the areas that matter
most to candidates and to Bupa: applying for a job, signing in, registering with a CV and finding the right job.
Business rules are checked quickly through the API, a small set of UI tests proves the real journeys work, and
the parts that need human judgement (the CAPTCHA on production, exploratory testing and language quality) stay
manual. Built this way, the suite gives fast, trustworthy feedback and stays cheap to maintain as the portal evolves.

\end{document}
